% Lösungen Einheit 3 von 4 – Klammern auflösen (zusammenbau v0.1)
\einheitenkopf{Einheit 3 von 4 · Klammern auflösen}

\erg{17}{a) ein Plus}
%% TODO Lösung der Erklärzeile 18f) fehlt in der Bank
\erg{18}{a) $4x + 3y - 7$ \quad b) $2a + 5b + 1$ \quad c) $8 + 6x - y$ \quad d) $7y + 2z + 9$ \quad e) $3b + 4c - 5$ \quad f) \ldots \quad g) $5x + 15$ \quad h) $5a - 2b - 4$ \quad i) $7a - 2b + 5 - c$ \quad j) $-2x - 10$ \quad k) $7x + 12$ \quad l) $4x + 23$}
\erg{19}{a) Das Minus vor der Klammer dreht beide Vorzeichen, aus $-5$ wird $+5$. Richtig: $12 - (a - 5) = 12 - a + 5 = 17 - a$. \quad b) Ja: $5 - (x - 2) = 5 - x + 2 = 7 - x$. \quad c) $5 \cdot (x + 2) = 5x + 10$}
